\defmodule {menumrg}

This is an interactive program to build, manipulate, compare, and
reduce lattice bases associated with LCGs or MRGs.
It uses the ``{\tt SI SI}'' representation.
The main purpose is to examine interactively what happens to a
basis when different operations are performed on its vectors.
This could be useful for ``debugging'' or to deal with special classes
of bases for which the standardized procedures may not work very well.
A variant of this module has been used, for instance, to obtain the
numerical results given in \cite{rCOU94a,rTEZ93a}.
In those cases, the shortest vectors and reduced bases were
computed in (partly) ``manual'' mode.

Two bases {\tt B1} and {\tt B2} and one vector {\tt Ve}
are created initially (automatically) and available to the user.
All along, {\tt B1} and {\tt B2} will be bases of the same lattice.
The program performs different operations, classified according to which
of those three objects they involve:
(a) the operations which affect only {\tt B1};
(b) those affecting {\tt B1} and {\tt B2};
(c) and those involving {\tt B1} and {\tt Ve}.
All those operations are available from a menu (and some submenus), as
shown in Figure~\ref{mainmenu}.
We now briefly explain the different menu options.

%%%%%%%%%%%%%%%%%%%%%%%%%%%%%%%%%%%%%%%%%%%%%%%%%%%%%%%
\setbox1=\vbox {\hsize = 6.0in
\code

  Operations on B1 :
    1. Building MRG basis
    2. Reading/Writing
    3. Duality verification
    4. Dualizing
    5. Square lengths
    6. Sorting
    7. Pre-reduction
    8. Shortest vector/Minkowski reduction
    9. Shortest/Minkowski verification

  Operations involving B1 and B2 :
    10. Copy B1 in B2
    11. Permute B1 and B2
    12. Equivalence of B1 and B2
    13. Relative coordinates
    14. Lexicographic order


  Operations involving B1 and Ve :
    15. Vector construction
    16. Relative coordinates
    17. Insertion of Ve in B1

\endcode
}

\begin{figure}[htb] \centering \boxit{\box1}
\caption { Main menu for the program {\tt menumrg}.}
\label {mainmenu}
\end{figure}


\begin {description}
%
\item [{\bf Operations on B1 :}]
\item [1. {\em Building MRG basis\/} :]
   Asks the user for a modulus, order, multiplier(s), and dimensions,
   then constructs the appropriate lattice basis in {\tt B1}.
\item [2. {\em Reading/Writing\/} :]
   Read/write the basis {\tt B1} from/to a file or the terminal.
\item [3. {\em Duality verification\/} :]
   Verifies whether the dual basis is ``correct'' by computing the
   appropriate scalar products.
\item [4. {\em Dualizing\/} :]
   Interchange the primal/dual bases.
\item [5. {\em Square lengths\/} :]
   Compute the squared vector lengths.
   The values can be displayed either in full (large integers)
   or in scientific notation.
\item [6. {\em Sorting\/} :]
   Sorts the vectors of {\tt B1} by increasing length.
\item [7. {\em Pre-reduction\/} :]
   Performs ``pairwise'' prereductions $T$ and $T^*$ of \cite{rDIE75a}
   in various combinations or orders.
   Can also ask for a given number of pre-reductions.
\item [8. {\em Shortest vector/Minkowski reduction\/} :]
   Computes the shortest vector in the lattice, or a MRLB,
   using the module {\tt REDBAS} (does not compute error bounds).
\item [9. {\em Shortest/Minkowski verification\/} :]
   Verifies whether the first vector of {\tt B1} is really a
   shortest vector, or whether the current basis is really a MRLB,
   by performing the branch-and-bound with error bounds
   (using either {\tt REDBAS2} or {\tt REDBAS3}).
   An option is also available to compute the squared distance between
   each vector of {\tt B1} and the subspace generated by all the other vectors,
   and to select the vectors for which that distance is smaller than
   the length of the (current) shortest vector in {\tt B1}.
%
\item [{\bf Operations involving B1 and B2 :}]
\item [10. {\em Copy B1 onto B2\/} :]
   Copies the basis {\tt B1} onto {\tt B2}.
\item [11. {\em Permute B1 and B2\/} :]
   Permutes {\tt B1} and {\tt B2}.
\item [12. {\em Equivalence of B1 and B2\/} :]
   Tells whether or not {\tt B1} and {\tt B2} generate the same lattice.
   Assumes that the dual basis is ``updated''.
\item [13. {\em Relative coordinates\/} :]
   Gives the relative coordinates of a vector of {\tt B1} relative to
   the basis {\tt B2}.
\item [14. {\em Lexicographic order\/} :]
   Displays the squared lengths of the vectors of {\tt B1} and {\tt B2}
   in lexicographic order.
   Assumes that the two bases are already sorted.

%
\item [{\bf Operations involving B1 and Ve :}]
\item [15. {\em Vector construction\/} :]
   Permits one to either initialize {\tt Ve} to zero; 
   add to {\tt Ve} a multiple of a vector of {\tt B1};
   increase the dimension of {\tt Ve} by one (the new coordinate is zero);
   shift the coordinates of {\tt Ve} to the right by one (provided that the
   last coordinate is zero); or perform a linear combination of iterates
   of the latter shifting operations (as in \cite{rCOU94a}).
\item [16. {\em Relative coordinates\/} :]
   Computes the coordinates of {\tt Ve} relative to {\tt B1}.
\item [17. {\em Insertion of Ve in B1\/} :]
   Inserts {\tt Ve} into the basis {\tt B1}, provided that {\tt Ve}
   belongs to the lattice generated by {\tt B1} and is primitive.
\end {description}
